\section{引言}
路径有序相位因子 $\text{P}\exp\left(ig\int A_{\mu} dx^{\mu}\right)$，即所谓威尔逊线\footnote{在本工作中，我们只关注非光锥威尔逊线与威尔逊圈；只要不产生歧义，就把它们简称为威尔逊线。}，是规范理论中的一个基本对象，它在把局域场连接起来以构造规范不变算符方面起着关键作用。对威尔逊线和威尔逊圈重正化的研究始于 Polyakov \cite{Polyakov:1980ca} 与 Gervais、Neveu \cite{Gervais:1979fv}。基于一维辅助场形式 \cite{Gervais:1979fv,Arefeva:1980zd, Arefeva:1980uy,Samuel:1978iy,Brandt:1978wc}，威尔逊线可乘重正化的全阶证明通过图方法在文献 \cite{Dotsenko:1979wb,Craigie:1980qs} 中、通过泛函方法在文献 \cite{Dorn:1986dt} 中给出；综述见文献 \cite{Dorn:1986dt}。

非光锥威尔逊线算符由通过类时或类空威尔逊线相连的局域算符构成，这保证了非局域算符的规范不变性。威尔逊线算符在多种语境中有广泛应用。例如，类空威尔逊线算符在确定准分布函数（quasi-PDF）\cite{Ji:2013dva} 与赝分布函数（pseudo-PDF）\cite{Radyushkin:2017cyf} 方面起着关键作用，而类时威尔逊线算符对于在势非相对论QCD（pNRQCD）有效理论框架内理解夸克偶素的衰变与产生至关重要 \cite{Pineda:1997bj,Brambilla:1999xf,Brambilla:2004jw,Brambilla:2001xy,Brambilla:2002nu,Brambilla:2020ojz,Brambilla:2021abf,Brambilla:2022rjd,Brambilla:2022ayc}。

近年来，人们对 quasi-PDF 和 pseudo-PDF 的研究兴趣日益增加；它们被定义为如下类型的威尔逊线算符的强子矩阵元：
\begin{eqnarray}
\label{eq:defqoperator}
\mathcal{O}_{\Gamma}(zv) &=& \bar{\psi}(zv)\Gamma W(zv,0)\psi(0),\\
\label{eq:defgoperator}
\mathcal{O}^{\mu\nu\alpha\beta} (zv)&=& g^2 F^{\mu\nu}(zv)W(zv,0)F^{\alpha\beta}(0),
\end{eqnarray}
其中 $\psi$ 是夸克场，$F^{\mu\nu}$ 是胶子场强张量，$\Gamma$ 表示一个狄拉克结构，$v$ 是一个四矢量，在闵氏时空中对类空（或类时）威尔逊线算符满足 $v^2=-1$（或 $v^2=1$）\footnote{在本节及后续讨论中，我们在闵氏时空中进行。然而，引入梯度流时需要过渡到欧氏时空，因为梯度流定义在欧氏时空中，这意味着对类空和类时威尔逊线算符都有 $v^2=1$。}，$z$ 是实数，而 $W(zv, 0)$ 是连接两个场的直线路径有序威尔逊线
\begin{eqnarray}
W(zv, 0) = \text{P}\exp\left(ig\int_0^z ds\, v\cdot A(sv) \right),
\end{eqnarray}
其中夸克情形取规范群的基础表示，胶子情形取伴随表示。不失一般性，我们在本工作中始终取 $z>0$。值得注意的是，基于短距离算符乘积展开（OPE）与大动量有效理论（LaMET）\cite{Braun:2007wv, Ji:2014gla,Izubuchi:2018srq}，quasi-PDF 和 pseudo-PDF 可以与标准光锥部分子分布函数（PDF）建立因子化关系。采用非光锥威尔逊线算符（而非类光算符）的强子矩阵元，其优点是非常适合格点计算。然而，这种便利的代价是非光锥威尔逊线算符高度非平凡的重正化。全面综述可参见文献 \cite{Ji:2020ect}。
此外，胶子型威尔逊线算符（例如 eq.~(\ref{eq:defgoperator}) 定义的算符）的真空期望值与 pNRQCD 框架下夸克偶素衰变与产生的研究以及 QCD 真空结构相关联（综述见文献 \cite{DiGiacomo:2000irz}）。

梯度流最初是作为在格点计算中正则化并压制紫外（UV）涨落的方法被引入的 \cite{Narayanan:2006rf,Luscher:2009eq}。它已被证明可用于格点 QCD 的标度设定以及局域矩阵元和关联函数的格点计算 \cite{Luscher:2010iy,Luscher:2011bx,Luscher:2013cpa,Luscher:2013vga,Borsanyi:2012zs,Sommer:2014mea,Suzuki:2013gza,Makino:2014taa,Harlander:2018zpi,FlavourLatticeAveragingGroupFLAG:2021npn,Leino:2021vop,Mayer-Steudte:2022uih,Leino:2022kgj}。除降噪效应外，梯度流还表现出一个显著性质：只要用重正化的物理参数和场来表达，正流时的复合算符即使在连续极限下仍保持有限 \cite{Luscher:2011bx}。这一性质使梯度流适合作为格点计算与微扰计算之间的匹配方案。值得指出的是，在小流时极限下梯度流不改变算符的红外行为，因此在从梯度流方案匹配到 $\overline{\rm MS}$ 类方案时，可以使用维数正则化来调节红外（IR）发散。

结合梯度流研究 quasi-PDF 的最初提议见文献 \cite{Monahan:2016bvm}。随后，针对用 eq.~(\ref{eq:defqoperator}) 中算符定义的、与夸克准PDF相关的零外动量强子矩阵元，文献 \cite{Monahan:2017hpu} 完成了包含完整流时依赖的单圈匹配计算。在本研究中，我们针对威尔逊线算符——如 eq.~(\ref{eq:defqoperator}) 与 eq.~(\ref{eq:defgoperator}) 所定义的那些——提出在小流时极限下从梯度流方案匹配到 $\overline{\rm MS}$ 方案的系统方法。我们的方法基于一维辅助场形式，并给出单圈阶匹配计算的完整细节。
微扰梯度流的近期应用及更多文献参见 refs.~\cite{Artz:2019bpr,Rizik:2020naq,Suzuki:2020zue,Harlander:2020duo,Boers:2020uqr,Shindler:2021bcx,Harlander:2022tgk,Mereghetti:2021nkt,Harlander:2022vgf,Brambilla:2021egm,Crosas:2023anw}。

本文其余内容安排如下。在第 \ref{sec:EFT} 节中，我们介绍一维辅助场形式，并详细阐述其在威尔逊线算符重正化中的应用。在第 \ref{sec:introGF} 节中，我们介绍梯度流形式，并对局域流算符实施小流时展开，得到匹配计算的核心公式。第 \ref{sec:oneloopcalc} 节专门计算至单圈阶的匹配系数。总结与结论见第 \ref{sec:summary} 节。在附录 \ref{app:quarkquasi} 中，我们给出与零外动量准PDF相联系的强子矩阵元的含完整流时依赖的详细匹配计算。我们将这些结果与文献 \cite{Monahan:2017hpu} 的已有结果以及第 \ref{sec:oneloopcalc} 节在小流时极限下的结果进行比较。
